\documentclass[11pt]{article}
\usepackage[a4paper,margin=24mm]{geometry}
\usepackage{amsmath,amssymb,booktabs,hyperref}
\hypersetup{colorlinks=true,urlcolor=blue,linkcolor=blue}
\newcommand{\dd}{\mathrm{d}}
\newcommand{\CM}{\mathcal{T}}
\title{Grouped recovery: conditional equivalence and projection limits}
\author{RM12 BP-R Stage 3 --- methodological review}
\date{7 October 2026}
\begin{document}
\maketitle
\noindent\textbf{Status.} M1 authorizes grouped-state qualification, not adoption
of the rules examined below. These are reviewer deductions from the supplied
Laval--Daganzo source and explicitly conditional algebra. They do not establish
a coupled ramp trajectory, control benefit, finite-summary impossibility or
useful gradient. No traffic simulation is used.

\section{State and sampled motion}
At a completed source sample, let $x$ be position (m), $v$ realized speed (m/s),
$\delta$ time step (s), $a_0$ zero-speed acceleration (m/s$^2$), $v_{\max}$ the
vehicle capability parameter (m/s), $g$ gravity and $G(x)$ road grade. Fixed lane
and type are implicit. The source Appendix A advances
\begin{align}
 x^+&=x+\delta v,\qquad
 w^+=v+\delta\{a_0(1-v/v_{\max})-gG(x)\}, \label{eq:sample}\\
 v^+&=\min\{w^+,U^+(x^+)\}. \label{eq:clip}
\end{align}
Here $w^+$ is next desired speed. $U^+$ denotes the consistently estimated
next downstream KW speed, from the model's own coupled field. Equation
\eqref{eq:clip} is not a substitute algorithm for determining that field.
At an intermediate phase $w^+$ cannot be discarded. Birth/node/crossing
existence and timing remain separate qualification obligations.

A lane-wise descriptor measure is
\begin{equation}
 \mu=\sum_{j=1}^{K}\delta_{\sigma_j},\qquad
 \sigma_j=(x_j,v_j,\theta_j,c_j),\qquad
 \mu^+=\CM_{s\,\#}\mu+\beta_s-\eta_s. \label{eq:measure}
\end{equation}
$\theta_j$ stores fixed motion/type and any required sample-phase data;
$c_j$ denotes a cumulative-count anchor (vehicle equivalents), if required by
the coupled representation, computed from its own traffic field at the declared
moving-boundary sample. It is not a personal vehicle identifier.
$s$ is the model's own traffic field, not reference truth; $\beta_s$ records
own-transfer births, and $\eta_s$ only justified removal. The pushforward maps
each signature by the same specified own-state update. This is a conditional
description, not a proof that the joined solver is complete. Descriptor count
is not physical traffic mass or additional TST population.

\section{When a group is lossless}
Group only identical \emph{full} signatures, retaining multiplicity $m_h$:
\begin{equation}
 \mu=\sum_{h=1}^{H}m_h\delta_{\sigma_h},\qquad
 \CM_{s\,\#}\mu=\sum_{h=1}^{H}m_h\delta_{\CM_s(\sigma_h)}.
 \label{eq:groups}
\end{equation}
This identity preserves the complete measure; $H\leq K$, with $H=K$ possible.
If the coupled solution is well defined and independent of descriptor names,
identical signatures see the same local field and have identical future
updates. Repeated identical constraint inequalities are redundant. Where a
qualified cumulative-count construction combines contributions by a minimum,
\begin{equation}
 \min\{C_0,C,C,\ldots\}=\min\{C_0,C,\ldots\}. \label{eq:idempotent}
\end{equation}
Thus multiplicity does not multiply the service penalty. The road field still
contains all continuous transferred mass. These arguments do not merge distinct
positions, speeds, count associations or types. Equal current speed alone is
insufficient. No assumption is made that sampled births share identical
placement, nor that a currently inactive constraint can be permanently removed.

\section{Joint location--speed association matters}
This is a local-input test of \eqref{eq:sample}--\eqref{eq:clip}, not a common-rate
traffic witness. On a level road, assign illustrative algebraic parameters
$\delta=1$ s, $a_0=3$ m/s$^2$, $v_{\max}=30$ m/s. Then the desired-speed map is
$\psi(v)=3+0.9v$. In two separated regions take the same prescribed local-input
speed estimates $U_L=15$ m/s and $U_R=20$ m/s. Positions $0$ m and $100$ m and
their one-step advances stay in their respective regions. Compare
\begin{equation}
 \mu_A=\delta_{(0,10)}+\delta_{(100,14)},\qquad
 \mu_B=\delta_{(0,14)}+\delta_{(100,10)}. \label{eq:swap}
\end{equation}
They share position and speed marginals, count, mean position and mean speed.
Both current speeds are below each assigned local cap. But
\begin{align}
 S_A^+&=\min(12,15)+\min(15.6,20)=27.6\ \mathrm{m/s},\\
 S_B^+&=\min(15.6,15)+\min(12,20)=27.0\ \mathrm{m/s}. \label{eq:sums}
\end{align}
Separate marginals cannot determine this speed-sum update. A coarse spatial bin
covering both regions can lose the same association, even with velocity bins.
Resolving the joint support may distinguish it. No claim is made that the
prescribed field is the eventual joined solution for both configurations, or
that this difference changes exit flow, TST or an optimal action.

\section{Count and mean position do not close histogram transport}
At one position update, use the source's current-speed advance $x^+=x+h$ with
$h=\delta v=2$ m and common $v=2$ m/s, $\delta=1$ s. All positions are initially
in the interval $[0,10)$ m:
\begin{equation}
 A=\{5.5,7.5\},\qquad B=\{4.5,8.5\}. \label{eq:positions}
\end{equation}
Both have count two, mean $6.5$ m and the same speed list. After translation,
$A^+=\{7.5,9.5\}$ stays inside while $B^+=\{6.5,10.5\}$ sends one record across
the boundary. Hence their projected next bin counts differ, despite equal
initial count/mean summaries. This does not defeat a summary retaining the
full subcell profile or a justified richer/restricted family.

For bins $B_a$ and projection $P_a\mu=\mu(B_a)$, the exact formula
\begin{equation}
 (P\CM_{s\,\#}\mu)_a=\int\mathbf{1}_{B_a}(\CM_s(z))\,\dd\mu(z)
 \label{eq:binupdate}
\end{equation}
requires the within-bin measure. Calling this a formula does not make it an
update from $P\mu$ alone. A declared reconstruction $R$ can instead define a
closed approximate map $P\CM_sR$, with births/coupling added consistently.
Its accuracy, activity and service are separate qualifications. A later bin
refinement cannot retrieve lost member locations from the reference.

\section{Conditional service and objective bounds}
For \emph{already qualified} trajectory contributions, suppose lower and upper
count contributions satisfy $L_h(z)\leq C_j(z)\leq U_h(z)$ for every member
$j$ of group $h$, on the same valid domain (including extended $+\infty$ values
outside domains of influence). Then
\begin{equation}
 \min\{C_0,\min_h L_h\}\ \leq\
 \min\{C_0,\min_j C_j\}\ \leq\
 \min\{C_0,\min_h U_h\}. \label{eq:envelope}
\end{equation}
This follows from monotonicity of minimum. It does not supply $L_h,U_h$ for
coupled unknown trajectories, sampled births or receiving interactions. Count
inequalities cannot simply be differentiated into pointwise flow inequalities.
If uncertainty straddles active/inactive configurations, a worst-case penalty
does not certify exact inactive behavior for the actual reference state.

Let $I(t)$ be cumulative offered population, including initial population, and
$E(t)$ cumulative final exits, in vehicle equivalents. With all entrance,
ramp and road waiting counted, system inventory is $I(t)-E(t)$. For common $I$,
finite-horizon total system time and its error satisfy
\begin{align}
 J_T&=\int_0^T [I(t)-E(t)]\,\dd t,\\
 |J_T-\widehat J_T|&\leq\int_0^T|E(t)-\widehat E(t)|\,\dd t.
 \label{eq:tst}
\end{align}
Units are vehicle-equivalent seconds. Full clearance requires certified tails
for both models. A genuine exit-count enclosure gives a corresponding TST
enclosure by integration; bin widths alone do not provide such an enclosure.
Even a small objective error does not prove correct or useful gradients.

\section*{Disposition}
Exact identical-signature grouping is a conditional equivalence rule. Broader
compact reconstruction remains unqualified. The sampled task, point merge,
priority, public data and TST are unchanged. The concrete grouping/approximation
decision S3-G2-Q-G is open; no candidate is adopted and Stage 4 is unopened.

\section*{Source locations}
Supplied Laval--Daganzo (2006), Section 2.2 and Appendix A, supplies the motion
and descriptor semantics. Daganzo--Laval RR-2003-2, Sections 2--4/Appendix, and
Simoni--Claudel, arXiv:1701.01026, Sections 2--3, support count/trajectory
coupling under their stated laws. Equations here are conditional reviewer
deductions, not a transferred theorem for the complete ramp task.
\end{document}
